\begin{equation}
    \begin{array}{lrll}
        \text{Formulas:} & \phi, \psi & ::= p \quad                     & \text{(Variabili proposizionali)} \\
                         &            & \mid \top^1 \quad               & \text{(Vero)}                     \\
                         &            & \mid \bot^0 \quad               & \text{(Falso)}                    \\
                         &            & \mid (\phi \land \psi)^{\times} & \text{(Congiunzione)}             \\
                         &            & \mid (\phi \lor \psi)^{+}       & \text{(Disgiunzione)}             \\
                         &            & \mid (\phi \to \psi)^{\to}      & \text{(Implicazione)}
    \end{array}
    \quad \text{con } \neg \phi \equiv \phi \to \bot^0
    \tag{IPC}
    \label{eq:ipc_grammar}
\end{equation}

\paragraph{Regole di formazione.}
Le regole di buona formazione stabiliscono i criteri grammaticali per determinare se una stringa di simboli costituisca un enunciato sensato del linguaggio, agendo come filtro per intercettare l'errore prima della fase di deduzione \cite{hindley1997}. 
Sotto la lente della Teoria Semplice dei Tipi, una formula è ben formata se e solo se la combinazione dei suoi sotto-enunciati rispetta rigorosamente l'arità dei costruttori di tipo ($\times, +, \to$) \cite{hindley1997}. Questo sbarramento sintattico 
assicura la coerenza strutturale della sintassi ($\mathcal{S}$), impedendo la formazione di espressioni autoreferenziali o degeneri.